% 自動生成（tools/mknext.py p3）。各章の道しるべ「この章の先に」から、後ろの巻の章を集めたもの（第13回）
\begin{longtblr}{colspec={Q[l,m,wd=58mm] X[l,m]}, width=\linewidth, rowhead=1, row{1}={bg=bkVol!14, font=\sffamily\bfseries\footnotesize},
  hlines={0.4pt,bkLine}, rowsep=1.5pt, cells={font=\footnotesize}}
この巻の章 & この道具を使う章 \\
{\bknextpill{波-1}\ ばねとおもりで考える：1自由度の振動} & {\bknextpill{電-9}\ スイッチを入れた瞬間：過渡現象} \\
{\bknextpill{波-7}\ 波が運ぶエネルギー：インピーダンスと反射} & {\bknextpill{電-13}\ 長い電線を読む：二端子対回路と分布定数回路} \\
{\bknextpill{波-12}\ 光の進み方：反射・屈折・レンズ} & {\bknextpill{電-15}\ 光も電波も同じ波：電磁波\\ \bknextpill{原-2}\ 光はつぶでもある：光電効果と光子} \\
{\bknextpill{波-13}\ 光の波：干渉・回折・偏光・散乱} & {\bknextpill{電-15}\ 光も電波も同じ波：電磁波\\ \bknextpill{原-2}\ 光はつぶでもある：光電効果と光子} \\
\end{longtblr}
